\documentclass[11pt]{article}
\usepackage[utf8]{inputenc}
\usepackage{amsmath,amssymb,amsthm}
\usepackage{hyperref}
\usepackage{mathtools}

\title{A Two-Parameter Linear Program for the Equation $x+y=mz$}
\author{}
\date{}

\newtheorem{proposition}{Proposition}
\newtheorem{remark}{Remark}

\begin{document}

\maketitle

\section{Motivation}

Let $p$ be prime and let $A \subseteq \mathbb{F}_p$. We want to bound the density of $A$
under the assumption that
\[
0 \notin A + A - mA,
\]
or equivalently that there are no triples $(x,y,z) \in A^3$ such that
\[
x+y=mz.
\]

The one-parameter LP from \texttt{search\_witness\_dil.py} and
\texttt{affine\_form\_lp.py} studies the joint distribution of indicators
$1_A(L(t))$ where $L$ is an affine form in a single variable $t$. This captures
the original Kravitz-set argument well, but it still leaves a large fake feasible
region for the $m$-sum-free problem.

The natural next step is to let the witness forms depend on two parameters.

\section{The Witness Forms}

Fix a finite family of affine forms
\[
L_r(u,v)=a_r u + b_r v + c_r \qquad (1 \le r \le n),
\]
where $a_r,b_r,c_r \in \mathbb{Q}$ and the denominators are assumed to be
invertible modulo $p$.

For each sign pattern $\varepsilon \in \{\pm 1\}^n$, define the atomic set
\[
S_{\varepsilon}
=
\bigcap_{r=1}^n \bigl( L_r^{-1}(A) \bigr)^{\varepsilon_r},
\]
where $X^{+1}=X$ and $X^{-1}=\mathbb{F}_p^2 \setminus X$.
Its density is
\[
a(\varepsilon)=\frac{|S_{\varepsilon}|}{p^2}.
\]

These are the LP variables. As usual they satisfy
\[
\sum_{\varepsilon} a(\varepsilon)=1,
\qquad
a(\varepsilon)\geq 0.
\]

If one of the forms, say $L_1$, is nonconstant, then $L_1(u,v)$ is uniformly
distributed on $\mathbb{F}_p$, so
\[
\delta(A)=\mathbb{P}(L_1(u,v)\in A)
=
\sum_{\varepsilon:\,\varepsilon_1=1} a(\varepsilon).
\]
The LP objective is to maximize this quantity.

\section{Forbidden Triples}

The key observation is that a triple of witness forms gives an empty atom whenever
it satisfies the target equation identically.

\begin{proposition}
Suppose that
\[
L_i(u,v)+L_j(u,v)-mL_k(u,v)\equiv 0
\]
as an identity in $(u,v)$. If $\varepsilon_i=\varepsilon_j=\varepsilon_k=1$,
then the atom $S_{\varepsilon}$ is empty.
\end{proposition}

\begin{proof}
Take $(u,v) \in S_{\varepsilon}$. Then $L_i(u,v),L_j(u,v),L_k(u,v)\in A$.
Since the displayed identity holds for every $(u,v)$, we get
\[
L_i(u,v)+L_j(u,v)=mL_k(u,v),
\]
which is forbidden by the assumption on $A$.
\end{proof}

Writing the identity coefficientwise, the condition is
\[
a_i+a_j-ma_k=0,\qquad
b_i+b_j-mb_k=0,\qquad
c_i+c_j-mc_k=0.
\]
Whenever this happens, the LP sets every atom with
$\varepsilon_i=\varepsilon_j=\varepsilon_k=1$ equal to $0$.

\section{Affine Symmetry on the Parameter Plane}

The analogue of translation invariance from the one-parameter LP is now invariance
under affine reparameterizations of $(u,v)$.

If
\[
T(u,v)=M\binom{u}{v}+t
\]
with $M \in \mathrm{GL}_2(\mathbb{F}_p)$ and $t \in \mathbb{F}_p^2$, then
\[
\bigl(1_A(L_r(T(u,v)))\bigr)_{r=1}^n
\]
has the same distribution as
\[
\bigl(1_A(L_r(u,v))\bigr)_{r=1}^n.
\]
Therefore two subfamilies of witness forms that differ by precomposition with an
affine automorphism of $\mathbb{F}_p^2$ must have the same intersection density.

The implementation in \texttt{two\_parameter\_lp.py} enforces this by assigning a
\emph{canonical shape} to every subset of the witness family.

\subsection*{Rank $1$ normalization}

Assume the coefficient vectors $(a_r,b_r)$ in the subset are all proportional.
Choose a nonconstant reference form
\[
L_{\mathrm{ref}}(u,v)=a_{\mathrm{ref}}u+b_{\mathrm{ref}}v+c_{\mathrm{ref}}.
\]
After an affine change of variables, we may normalize this reference form to $u$.
Every other form in the subset becomes
\[
\lambda u + \bigl(c-\lambda c_{\mathrm{ref}}\bigr),
\]
where $(a,b)=\lambda(a_{\mathrm{ref}},b_{\mathrm{ref}})$.

So the normalized data are the triples
\[
(\lambda,0,c-\lambda c_{\mathrm{ref}}).
\]
The code tries all possible reference forms and keeps the lexicographically smallest
normalized tuple.

\subsection*{Rank $2$ normalization}

Now assume the subset contains two forms with linearly independent coefficient
vectors. Choose an ordered pair $(L_i,L_j)$ of such forms. After an affine change
of variables, one can normalize them to
\[
L_i(u,v)=u,\qquad L_j(u,v)=v.
\]
Every other form in the subset is then expressed in these normalized coordinates.

Concretely, if
\[
L_i=a_i u+b_i v+c_i,\qquad L_j=a_j u+b_j v+c_j,
\]
and the coefficient vectors are independent, then the normalization is obtained by
solving:
\[
\left(M^T\right)
\begin{pmatrix}
a_i & a_j \\
b_i & b_j
\end{pmatrix}
= I_2,
\qquad
\begin{pmatrix}
a_i & b_i \\
a_j & b_j
\end{pmatrix}
t
=
\begin{pmatrix}
-c_i\\
-c_j
\end{pmatrix}.
\]
The code applies this normalization to every ordered independent pair and again
keeps the lexicographically smallest normalized tuple.

\section{The Resulting LP}

For a fixed witness family $\mathcal{L}=\{L_1,\dots,L_n\}$, the two-parameter LP is:

\begin{itemize}
\item Variables: the atom densities $a(\varepsilon)$.
\item Positivity and mass:
\[
a(\varepsilon)\geq 0,\qquad \sum_{\varepsilon} a(\varepsilon)=1.
\]
\item Forbidden-triple constraints: if $L_i+L_j-mL_k\equiv 0$, then every atom
with $\varepsilon_i=\varepsilon_j=\varepsilon_k=1$ has density $0$.
\item Affine-shape constraints: any two subsets of witness forms with the same
canonical shape have equal aggregate density.
\item Objective: maximize
\[
\sum_{\varepsilon:\,\varepsilon_{r_0}=1} a(\varepsilon),
\]
where $L_{r_0}$ is any nonconstant witness form.
\end{itemize}

\section{Implementation Notes}

The file \texttt{two\_parameter\_lp.py} implements the model.

\begin{itemize}
\item \texttt{linear\_form\_2d(a,b,c,label)} stores the form $au+bv+c$.
\item \texttt{canonical\_subset\_shape} computes the affine normal form of a subset.
\item \texttt{build\_two\_parameter\_lp(m, forms)} constructs the LP.
\item \texttt{evaluate\_two\_parameter\_bound(m, forms)} solves it and returns the
resulting upper bound for $\delta(A)$.
\end{itemize}

Two simple witness families already illustrate the behavior:

\begin{itemize}
\item Using only
\[
u,\quad v,\quad \frac{u+v}{m},
\]
the LP gives the bound $\delta(A)\leq 2/3$.
\item Adding the extra direction
\[
\frac{u-v}{m}
\]
improves the bound to $\delta(A)\leq 4/7$ in the initial experiments.
\end{itemize}

These are not yet competitive with the existing best bounds, but they show that the
two-parameter model behaves differently from the one-parameter LP and can encode
new structural information.

\end{document}
